\chapter{Introduction}
\label{ch:introduction}

% ─────────────────────────────────────────────────────────────────────────────
\section{What This Manual Covers}
% ─────────────────────────────────────────────────────────────────────────────

This manual is the authoritative reference for the \texttt{gert.ops} Domain Kit---a specialized vocabulary layer for authoring operations runbooks within the gert v2 framework. The \texttt{gert.ops} Kit implements the DRI (Directly Responsible Individual) accountability model, providing step types, governance primitives, and evidence-capture patterns designed for SRE and DevOps teams running production operations, change management, and incident response.

The manual covers:

\begin{itemize}
  \item The DRI accountability model and role taxonomy
  \item \texttt{gert.ops} step types: \texttt{ops.cli}, \texttt{ops.manual}, \texttt{ops.approval}, \texttt{ops.change-request}, \texttt{ops.incident}
  \item Role-based gating, approval workflows, and SLA timeout policies
  \item Evidence types: screenshots, command output, attestations with SHA256 verification
  \item Change request workflows: pre-execution approval, evidence capture, rollback injection
  \item Incident response workflows: declaration, triage, escalation, resolution
  \item Migration from v1 DRI runbooks to gert v2 + \texttt{gert.ops}
  \item Complete field reference for all \texttt{gert.ops} schema elements
\end{itemize}

This is \emph{not} a general introduction to gert v2 or Domain Kits. For core gert concepts (runbook structure, core step types, the Domain Kit model), see the \emph{gert v2 Design Document}. For guidance on developing your own Domain Kit, see the \emph{Domain Kit Development Guide}.

% ─────────────────────────────────────────────────────────────────────────────
\section{Audience and Prerequisites}
% ─────────────────────────────────────────────────────────────────────────────

This manual is written for:

\begin{description}
  \item[SRE and DevOps Engineers] — Writing runbooks for deployments, rollbacks, infrastructure changes, and service maintenance.
  \item[Operations Leads] — Designing change management processes, approval matrices, and escalation policies.
  \item[Incident Commanders] — Authoring incident response playbooks with triage steps and evidence requirements.
  \item[Change Managers] — Defining governance policies, approval gates, and audit trail requirements for compliance.
\end{description}

\paragraph{Prerequisites.}
You should be familiar with:

\begin{itemize}
  \item \textbf{gert v2 core concepts} — runbook structure, core step types (\texttt{cli}, \texttt{manual}, \texttt{tool}, \texttt{branch}, \texttt{iterate}), variables, and the execution model. See \emph{gert v2 Design Document}, §1--§3.
  \item \textbf{Domain Kit concepts} — what a Domain Kit is, how Kit-specific step types compile (lower) into core primitives, and the relationship between authoring vocabulary and runtime execution. See \emph{gert v2 Design Document}, §4.
  \item \textbf{YAML syntax} — \texttt{gert.ops} runbooks are authored in YAML. You do not need to know JSON Schema or Go.
\end{itemize}

If you are new to gert, start with the core documentation before reading this manual. The \texttt{gert.ops} Kit builds on core concepts and assumes you understand the foundation.

% ─────────────────────────────────────────────────────────────────────────────
\section{What is DRI?}
% ─────────────────────────────────────────────────────────────────────────────

\textbf{DRI} (Directly Responsible Individual) is an operational accountability model that designates exactly one person as accountable for a task, change, or incident at any given moment. The DRI is not necessarily the person \emph{doing} the work---they are the person \emph{responsible} for the outcome.

\paragraph{Origin.}
The DRI model originated at Apple and has been widely adopted in SRE and DevOps cultures. In operational contexts, DRI provides clarity during high-stress events: when an incident occurs or a change is being deployed, everyone knows who has authority to make decisions and who is accountable if things go wrong.

\paragraph{Why DRI Matters.}
Operations teams often face situations where multiple people are involved (a deployment might require an engineer, a database admin, and a manager's approval), but no one is clearly accountable. This creates confusion:
\begin{itemize}
  \item Who decides whether to proceed with a risky step?
  \item Who is authorized to approve a rollback?
  \item Who escalates if a timeout is reached?
  \item Who signs off on the evidence after completion?
\end{itemize}

The DRI model resolves this by making accountability explicit. At every moment during a runbook execution, there is one DRI. If the DRI is unavailable, responsibility transfers to a designated backup via the escalation chain.

\paragraph{DRI in gert.ops.}
The \texttt{gert.ops} Domain Kit implements DRI accountability through:
\begin{itemize}
  \item A role taxonomy: \texttt{dri}, \texttt{approver}, \texttt{change-manager}, \texttt{incident-commander}, \texttt{responder}, \texttt{observer}
  \item Step-level role requirements: \texttt{requires\_role} field enforces that only authorized roles can execute a step
  \item Approval gates: \texttt{ops.approval} step pauses execution until specified roles sign off
  \item Evidence signatures: every evidence item records which role captured it
  \item SLA enforcement: timeouts trigger escalation to backup roles if the primary DRI is unresponsive
\end{itemize}

% ─────────────────────────────────────────────────────────────────────────────
\section{How gert.ops Implements DRI}
% ─────────────────────────────────────────────────────────────────────────────

The \texttt{gert.ops} Kit extends gert v2 core with five domain-specific step types and additional metadata fields. These additions compile down to core primitives at build time, so the gert runtime never sees \texttt{gert.ops}-specific vocabulary. This compilation model ensures that the core engine remains stable and domain-agnostic.

\subsection{Step Types}

\begin{description}
  \item[\texttt{ops.cli}] — Extends core \texttt{cli} with audit logging, command allowlist enforcement, and environment variable blocking.
  \item[\texttt{ops.manual}] — Extends core \texttt{manual} with role requirements, SLA timeouts, and evidence type constraints.
  \item[\texttt{ops.approval}] — An approval gate: pauses execution until specified roles sign off. Compiles into a \texttt{manual} step with approval-check logic.
  \item[\texttt{ops.change-request}] — Wraps a sequence of steps in a change management context, recording change ID, approval timestamp, and automatic rollback step injection on failure.
  \item[\texttt{ops.incident}] — Wraps a sequence of steps in an incident response context, assigning incident ID, incident-commander role, and escalation chain.
\end{description}

\subsection{Metadata Fields}

\texttt{gert.ops} runbooks include additional metadata:

\begin{description}
  \item[\texttt{meta.roles}] — A map of role names to identities (person, team, or \texttt{@on-call}). Example: \texttt{dri: alice@example.com}, \texttt{approver: ops-team}.
  \item[\texttt{meta.change-id}] — Optional reference to an external change management system (e.g., JIRA ticket, ServiceNow change record).
  \item[\texttt{meta.sla}] — Execution window and escalation policy. Defines the maximum allowed duration and who to notify on timeout.
\end{description}

\subsection{Evidence Types}

\texttt{gert.ops} defines three evidence types that extend core evidence capture:

\begin{description}
  \item[\texttt{ops.evidence.screenshot}] — An image attachment with SHA256 hash, MIME type, and role signature.
  \item[\texttt{ops.evidence.command-output}] — Captured stdout/stderr from a \texttt{cli} step, with SHA256 and role signature.
  \item[\texttt{ops.evidence.attestation}] — Free-text evidence with timestamp and role signature (e.g., ``I verified the database replica lag is under 100ms'').
\end{description}

\subsection{Governance Integration}

\texttt{gert.ops} wires approval gates into gert core's governance hooks. When a \texttt{ops.approval} gate is reached:

\begin{enumerate}
  \item The runtime pauses and emits a \texttt{governance/approvalRequired} trace event.
  \item The runtime waits for approval decisions to be submitted via the \texttt{gert exec approve} command or API.
  \item Once the minimum number of approvals from specified roles is met, execution resumes.
  \item If the SLA timeout elapses before approvals are received, the runtime triggers the escalation policy (either fail the run or escalate to a backup approver).
\end{enumerate}

% ─────────────────────────────────────────────────────────────────────────────
\section{Document Structure}
% ─────────────────────────────────────────────────────────────────────────────

The remaining chapters are organized as follows:

\begin{description}
  \item[Chapter 1: DRI Concepts] — The DRI accountability model, role taxonomy, escalation chains, and SLA policies.
  \item[Chapter 2: Schema Reference] — The \texttt{gert.ops} Kit schema, step types, and metadata fields.
  \item[Chapter 3: Authoring Guide] — How to write a \texttt{gert.ops} runbook from scratch, with complete examples.
  \item[Chapter 4: Governance] — Role-based gating, allowlist enforcement, environment variable blocking, redaction, and the audit trail.
  \item[Chapter 5: Evidence and Tracing] — The evidence model, timeline projection, SHA256 verification, and compliance bundles.
  \item[Chapter 6: Change Request Workflow] — End-to-end walkthrough of a change request from pre-approval to closure.
  \item[Chapter 7: Incident Response Workflow] — End-to-end walkthrough of an incident from declaration to resolution.
  \item[Chapter 8: Migration] — Migrating v1 DRI runbooks to gert v2 + \texttt{gert.ops}.
  \item[Chapter 9: Reference] — Complete field reference tables for all \texttt{gert.ops} schema elements.
\end{description}

Each chapter is designed to be self-contained. If you are writing a deployment runbook, you can start with Chapter 3 (Authoring Guide) and refer to Chapter 2 (Schema Reference) and Chapter 9 (Field Reference) as needed. If you are designing an incident response playbook, start with Chapter 7. If you are migrating existing runbooks, start with Chapter 8.
